\section{Combat}

\subsection{Dmg-Formular}
\label{subsec:dmg-formula}

\subsubsection{Top-Level Formula}

The total damage of a single hit is defined as:

\[
\text{FinalDMG} = (\text{BaseEffektiveDMG} + \text{IncreasedDMG}) \times \text{MoreDMG}
\]

with:

\begin{itemize}
	\item \textbf{BaseEffektiveDMG} – base damage after all additive composition and global stat scaling.
	\item \textbf{IncreasedDMG} – all applicable \emph{increased} damage modifiers (additive layer).
	\item \textbf{MoreDMG} – all applicable \emph{more} damage modifiers (multiplicative layer).
\end{itemize}

% ---------------------------------------------------------------------------
\subsubsection{Main Components}

\paragraph{BaseEffektiveDMG}

\[
\text{BaseEffektiveDMG} = \text{BaseDMG} \times \text{DMGEffektModifier}
\]

\begin{itemize}
	\item \textbf{BaseDMG} is the total raw damage of the hit after base values, added damage, conversion and gain have been resolved (per damage type, then summed).
	\item \textbf{DMGEffektModifier} is a global stat-based multiplier derived from the attacker’s attributes.
\end{itemize}

\paragraph{IncreasedDMG}

Conceptually:

\[
\text{IncreasedDMG} = \text{BaseEffektiveDMG} \times \text{dmgIncreaseBuckets}
\]

Internally, this is calculated per damage type and then aggregated:

\[
\text{IncreasedDMG} = \sum_T \left(\text{BaseEffektiveDMG}_T \times \text{incSum}_T\right)
\]

where:

\begin{itemize}
	\item $T$ iterates over all damage types.
	\item $\text{incSum}_T$ is the sum of all applicable \emph{increased} modifiers for type $T$.
	\item \texttt{dmgIncreaseBuckets} is the effective combined increase when viewed on aggregate level.
\end{itemize}

\paragraph{MoreDMG}

\[
\text{MoreDMG} = \prod_i \text{moreFactor}_i
\]

All \emph{more} modifiers are collected as multiplicative factors and multiplied together. Many of these are conditional (e.g. ``more damage while standing'', ``more damage vs frozen enemies'').

% ---------------------------------------------------------------------------
\subsubsection{BaseDMG Breakdown}

The total base damage of a hit is defined as:

\[
\text{BaseDMG} = \text{SkillBaseDMG} + \text{AddedDMG} + \text{ConvertedDMG} + \text{GainedDMG}
\]

Internally, all damage is tracked per damage type $T$ (Physical, Fire, Cold, Lightning, Chaos, \dots). For each type:

\begin{itemize}
	\item \textbf{SkillBaseDMG} – base damage values of the skill per type.
	\item \textbf{AddedDMG} – flat added damage per type (e.g. ``+20 Cold Damage'').
	\item \textbf{ConvertedDMG} – damage that is moved from one type to another (no duplication).
	\item \textbf{GainedDMG} – damage that is gained as \emph{extra} from another type (duplication, ``gain X\% of A as B'').
\end{itemize}

Aggregated:

\[
\text{BaseDMG} = \sum_T \text{BaseDMG}_T
\]

% ---------------------------------------------------------------------------
\subsubsection{DMGEffektModifier and StatModifier}

\[
\text{DMGEffektModifier} = \text{StatModifier}
\]

\[
\text{StatModifier} = f(\text{BaseAttribute} + \text{GearModifier} + \text{OtherModifier})
\]

\begin{itemize}
	\item \textbf{BaseAttribute} – base attribute values from level progression.
	\item \textbf{GearModifier} – attribute bonuses from equipment.
	\item \textbf{OtherModifier} – other attribute sources (buffs, passives, etc.).
	\item $f(\cdot)$ – a defined scaling function that maps the final attribute value to a global damage multiplier
	(e.g. INT = 100 $\Rightarrow$ \texttt{StatModifier} = 2.5 = 250\%).
\end{itemize}

This modifier is \emph{global} and applies to the entire damage of the hit, independent of damage type.

% ---------------------------------------------------------------------------
\subsubsection{dmgIncreaseBuckets (Increased Layer)}

\texttt{dmgIncreaseBuckets} represents the combined effect of all \emph{increased} damage modifiers. Internally, this is resolved per damage type $T$:

\[
\text{incSum}_T = \sum_j \text{increasedModifier}_{T,j}
\]

where the sum includes:

\begin{itemize}
	\item Global increased modifiers (e.g. ``increased Damage'').
	\item Type-based modifiers (e.g. ``increased Physical Damage'', ``increased Cold Damage'', ``increased Elemental Damage'').
	\item Tag-based modifiers (e.g. ``increased Spell Damage'', ``increased Attack Damage'',
	``increased Projectile Damage'', ``increased AoE Damage''), if the skill has the corresponding tag.
	\item Conditional increased modifiers (e.g. ``increased Damage while on Full Life'', ``increased Damage vs Frozen''),
	if their condition is satisfied.
\end{itemize}

Application per type:

\[
\text{FinalIncLayer}_T = \text{BaseEffektiveDMG}_T \times (1 + \text{incSum}_T)
\]

Aggregated on hit level:

\[
\text{IncreasedDMG} = \sum_T \left(\text{BaseEffektiveDMG}_T \times \text{incSum}_T\right)
\]

and

\[
\text{preMoreTotal} = \text{BaseEffektiveDMG} + \text{IncreasedDMG}
= \sum_T \text{FinalIncLayer}_T
\]

% ---------------------------------------------------------------------------
\subsubsection{MoreDMG (Multiplicative Layer)}

All \emph{more} modifiers form their own multiplicative layer:

\[
\text{MoreDMG} = \prod_i \text{moreFactor}_i
\]

Examples:

\begin{itemize}
	\item ``20\% more Damage while standing'' $\rightarrow$
	factor 1.2 if the attacker is standing, otherwise 1.0.
	\item ``50\% more Damage'' $\rightarrow$ factor 1.5.
	\item ``X\% more Damage vs Frozen Enemies'' $\rightarrow$
	factor applied only if the target is frozen.
\end{itemize}

The final player damage (before defenses) is:

\[
\text{FinalDMG\_beforeDef} = \text{preMoreTotal} \times \text{MoreDMG}
\]

which matches the top-level formula:

\[
\text{FinalDMG} = (\text{BaseEffektiveDMG} + \text{IncreasedDMG}) \times \text{MoreDMG}
\]

% ---------------------------------------------------------------------------
\subsubsection{Step-by-Step Pipeline}

The following describes the detailed pipeline per hit.

\paragraph{Step 0: Build HitContext}

Before any numerical damage computation, a \textbf{HitContext} is constructed and frozen for this hit. It contains:

\begin{itemize}
	\item The skill and its tags (e.g. \texttt{attack}, \texttt{spell}, \texttt{projectile},
	\texttt{cold}, \texttt{elemental}, \texttt{aoe}, \dots).
	\item Attacker attributes (STR, DEX, CON, INT, WIL).
	\item Effects from gear, passives, buffs and debuffs.
	\item State flags (standing / moving, on full life, target frozen, etc.).
	\item Optionally critical hit information (depending on system design).
\end{itemize}

All conditional modifiers read \emph{only} from this HitContext, never from intermediate values of the damage calculation.

% ---------------------------------------------------------------------------
\paragraph{Step 1: Build BaseDMG per Damage Type}

\begin{enumerate}
	\item \textbf{Initialize SkillBaseDMG}
	
	\[
	base_T = \text{SkillBaseDMG}_T
	\]
	
	\item \textbf{Apply AddedDMG}
	
	\[
	base_T \mathrel{+}= \sum \text{AddedDMG}_T
	\]
	
	\item \textbf{Apply Conversions (with clamping as failsafe)}
	
	Design rule: normally, only a single conversion source per damage type is used in content.
	The engine supports multiple conversions per source type $S$ and normalizes them to a maximum of 100\%.
	
	For each source type $S$:
	
	\begin{enumerate}
		\item Store original value: \quad $\text{sourceBase} = base_S$
		\item Collect all conversion percentages $c_i$ for $S \rightarrow T_i$
		\item Compute
		\[
		\text{totalRaw} = \sum_i c_i
		\]
		\item If $\text{totalRaw} \leq 1.0$:
		\[
		\text{effConv}_i = c_i
		\]
		\item If $\text{totalRaw} > 1.0$ (failsafe):
		\[
		\text{scale} = \frac{1.0}{\text{totalRaw}}, \quad
		\text{effConv}_i = c_i \times \text{scale}
		\]
		\item Apply conversion:
		\[
		base_S = \text{sourceBase} \times (1 - \sum_i \text{effConv}_i)
		\]
		\[
		base_{T_i} \mathrel{+}= \text{sourceBase} \times \text{effConv}_i
		\]
	\end{enumerate}
	
	Converted damage is not converted again in V1 (no chaining like Physical $\rightarrow$ Cold $\rightarrow$ Lightning).
	
	\item \textbf{Apply GainedDMG (``Gain as Extra'')}
	
	After all conversions, gains are applied based on the already converted values:
	
	For each gain effect ``Gain $g$ of A as B'':
	
	\[
	\text{extra}_B \mathrel{+}= base_A \times g
	\]
	
	Then:
	
	\[
	base_T \mathrel{+}= \text{extra}_T
	\]
	
	Gains duplicate damage and may freely stack (no normalization).
\end{enumerate}

After Step~1:

\[
\text{BaseDMG}_T = base_T
\quad\text{and}\quad
\text{BaseDMG} = \sum_T \text{BaseDMG}_T
\]

% ---------------------------------------------------------------------------
\paragraph{Step 2: Apply StatModifier (DMGEffektModifier)}

Compute the global stat-based multiplier:

\[
\text{StatModifier} = f(\text{BaseAttribute} + \text{GearModifier} + \text{OtherModifier})
\]

Set:

\[
\text{DMGEffektModifier} = \text{StatModifier}
\]

Apply per type:

\[
\text{BaseEffektiveDMG}_T = \text{BaseDMG}_T \times \text{DMGEffektModifier}
\]

Aggregate:

\[
\text{BaseEffektiveDMG} = \sum_T \text{BaseEffektiveDMG}_T
\]

% ---------------------------------------------------------------------------
\paragraph{Step 3: Compute IncreasedDMG (dmgIncreaseBuckets)}

For each damage type $T$:

\begin{enumerate}
	\item Collect all applicable \emph{increased} modifiers from the HitContext
	(global, type-based, tag-based, conditional).
	\item Sum them:
	\[
	\text{incSum}_T = \sum_j \text{increasedModifier}_{T,j}
	\]
	\item Apply:
	\[
	\text{FinalIncLayer}_T = \text{BaseEffektiveDMG}_T \times (1 + \text{incSum}_T)
	\]
\end{enumerate}

On aggregate:

\[
\text{IncreasedDMG} = \sum_T \left(\text{BaseEffektiveDMG}_T \times \text{incSum}_T\right)
\]

\[
\text{preMoreTotal} = \text{BaseEffektiveDMG} + \text{IncreasedDMG}
= \sum_T \text{FinalIncLayer}_T
\]

% ---------------------------------------------------------------------------
\paragraph{Step 4: Compute MoreDMG}

Initialize:

\[
\text{MoreDMG} = 1.0
\]

For each applicable \emph{more} modifier in the HitContext:

\[
\text{MoreDMG} \mathrel{*}= \text{moreFactor}_i
\]

Finally:

\[
\text{FinalDMG\_beforeDef} = \text{preMoreTotal} \times \text{MoreDMG}
\]

which matches the top-level definition:

\[
\text{FinalDMG} = (\text{BaseEffektiveDMG} + \text{IncreasedDMG}) \times \text{MoreDMG}
\]

% ---------------------------------------------------------------------------
\paragraph{Step 5: (Optional) Enemy Defenses}

Enemy defenses (armor, resistances, block, dodge, etc.) are applied after the player-side
damage pipeline and can be specified separately. Typically:

\begin{itemize}
	\item Resistances or mitigation per damage type are applied to \texttt{FinalDMG\_beforeDef}.
	\item Additional mechanics (block, dodge, avoidance) are resolved either in the hit resolution
	before this pipeline or as a final reduction step.
\end{itemize}

% ---------------------------------------------------------------------------
\paragraph{Order in Code \& Role of HitContext}

In code, the damage pipeline should follow the exact order below:

\begin{enumerate}
	\item Build and freeze the \textbf{HitContext} (skill, tags, attributes, buffs, states, etc.).
	\item Resolve \textbf{BaseDMG} per damage type:
	base values, added damage, conversions (with clamping), gains.
	\item Apply \textbf{StatModifier} to obtain \textbf{BaseEffektiveDMG} per type.
	\item Compute \textbf{IncreasedDMG} per type and aggregate to \texttt{preMoreTotal}.
	\item Collect and apply all \textbf{MoreDMG} factors to obtain \texttt{FinalDMG\_beforeDef}.
	\item Apply enemy defenses (optional).
\end{enumerate}

The HitContext is the single source of truth for all conditions. No modifier may depend
on intermediate numerical results of this pipeline. This keeps the order of operations
inside each layer irrelevant (additive increased, multiplicative more) and makes the
damage calculation deterministic and debuggable.
